\section{The \texorpdfstring{$\lambprot$}{lambda coin} calculus}~\label{sec:calculus}
In this section we present 
$\lambprot$ (read ``lambda coin''), which is the
simply typed lambda calculus extended with booleans ($1$ and $0$), an if-then-else construction, and a coin $\coin$.
%
Terms are inductively defined by 
\[
  t := x\mid \lambda x.t \mid tt \mid 1\mid 0\mid \ite ttt\mid \coin
\]

The rewrite system is given in Table~\ref{tab:TRS}. The rules $t\topr[p]r$ mean that $t$ reduces with probability $p$ in one step to $r$, where the sum of probabilities of reducing one redex is $1$. In particular, every non-contextual rule has probability $1$, since there is only one rule per redex, except for the coin, which reduces with probability $\frac 12$ to $0$ and probability $\frac 12$ to $1$. If $t\topr[p_1]r_1\dots\topr[p_n]r_n$ we may write $t\topr^*r_n$, where $p=\prod_{i=1}^np_i$.

The reduction rules are intentionally as permissive as possible (even the branches of an if-then-else can be reduced) in order to analyse its (lack of) confluence. The type system is given in Table~\ref{tab:TS}.

\begin{table}[t]
  \[
    \begin{array}{r@{\ }l@{\qquad}r@{\ }l@{\qquad}r@{\ }l}
      (\lambda x.t)r &\topr t[r/x]
      &
      \ite 1tu &\topr t
      &
      \coin &\topr[{\frac 12}] 1
      \\
      &&
      \ite 0tu &\topr u
      &
      \coin &\topr[{\frac 12}] 0
    \end{array}
  \]
  \smallskip
  \[
  \begin{array}{c c c}
    \infer{\lambda x.t\topr[p]\lambda x.r}{t\topr[p] r} & \infer{ts\topr[p] rs}{t\topr[p] r} & \infer{st\topr[p] sr}{t\topr[p] r}\\[5pt]
    \multicolumn{3}{c}{\infer{\ite{t}{s_1}{s_2}\topr[p]\ite{r}{s_1}{s_2}}{t\topr[p] r}}\\[5pt]
    \multicolumn{3}{c}{\infer{\ite{t}{s_1}{s_2}\topr[p]\ite{t}{r_1}{s_2}}{s_1\topr[p] r_1}}\\[5pt]
    \multicolumn{3}{c}{\infer{\ite{t}{s_1}{s_2}\topr[p]\ite{t}{r_1}{r_2}}{s_2\topr[p] r_2}}
  \end{array}
  \]
  \caption{Rewrite system for $\lambprot$.}
  \label{tab:TRS}
\end{table}

\begin{table}[t]
  \[
    A:= \mathbb{B}  \mid A\rightarrow A
  \]
  \[
    \infer[\mathsf{ax}]{\Gamma,x:A\vdash x:A}{} 
    \qquad
    \infer[\mathsf{ax}_1]{\Gamma\vdash 1:\mathbb{B}}{}
    \qquad
    \infer[\mathsf{ax}_0]{\Gamma\vdash 0:\mathbb{B}}{}
  \]
  \[
    \infer[\rightarrow_i]{\Gamma\vdash\lambda x.t:A\rightarrow B}
    	{\Gamma,x:A\vdash t:B}
	  \qquad 
    \infer[\rightarrow_e]{\Gamma\vdash tr:B}
    	{\Gamma\vdash t:A\rightarrow B & \Gamma\vdash r:A}
  \]
  \[
    \infer[\mathsf{ax}_{\coin}]{\Gamma\vdash \coin:\mathbb{B}}{} 
    \qquad
    \infer[\mathsf{if}]{\Gamma\vdash\ite t u v:A}{\Gamma\vdash t:\mathbb{B}& \Gamma\vdash u:A & \Gamma\vdash v:A}
  \]
  \caption{Type system for $\lambprot$.}
  \label{tab:TS}
\end{table}

Strong normalization for this calculus follows trivially from Tait's proof for the simply typed $\lambda$-calculus, extended with booleans. The only reduction added is the probabilistic coin toss and it takes at most one step for each operator. Hence, using the rewriting over probabilistic distributions techniques from~\cite{DCM17}, we only need to show local confluence in order to achieve global confluence altogether. This is an adaptation of Newman's lemma for probabilistic calculi (see~\cite{DCM17} for a longer discussion about probabilistic confluence).

Clearly, $\lambprot$ is not confluent, as already seen in the preliminaries example (see Figure~\ref{fig:noconfl}). It is easy to see that these distributions represent different results.

The divergence stems from three characteristics of the calculus:
\begin{enumerate}
  \item Lack of a reduction strategy.
  \item Probabilistic reductions.
  \item Duplication of variables.
\end{enumerate}
  
Removing just one of these elements renders the system confluent, however each modification comes with its own trade-off. We will examine each case, one by one, in the following section.
